\exercisesection{}

\begin{enumerate}
\def\labelenumi{\arabic{enumi})}
\item
  Let \(G\) be a compact group. Show that every continuous representation \(\varphi\) of \(G\) in \(\mathbf{R}_+^*\) is such that \(\varphi(G) = \{1\}\) . From this, deduce that a relatively invariant positive measure on \(G\) is invariant.
\item
  Let G be a topological group, such that the commutator subgroup of G is dense in G. Show that every continuous representation \(\varphi\) of G in a Hausdorff abelian group satisfies \(\varphi(G) = \{e\}\) . From this, deduce that if G is locally compact, then every relatively invariant complex measure on G is invariant. In particular, G is unimodular.
\item
  Let \(G\) be a locally compact group, \(\mu\) a left Haar measure on \(G\) , and \(\nu = \Delta_G^{-1/2} \cdot \mu\) . Show that
\end{enumerate}

\[
\boldsymbol {\gamma} (s) \nu = \Delta_ {\mathrm{G}} (s) ^ {1 / 2} \nu , \quad \boldsymbol {\delta} (s) \nu = \Delta_ {\mathrm{G}} (s) ^ {1 / 2} \nu , \quad \stackrel {\vee} {\nu} = \nu .
\]

\begin{enumerate}
\def\labelenumi{\arabic{enumi})}
\setcounter{enumi}{3}
\item
  Let \(G, G'\) be two locally compact groups, \(V\) (resp. \(V'\) ) an open neighborhood of the neutral element of \(G\) (resp. \(G'\) ), \(\varphi\) a local isomorphism of \(G'\) with \(G\) , defined on \(V'\) , such that \(\varphi(V') = V\) . Show that \(\Delta_G \circ \varphi\) is the restriction of \(\Delta_{G'}\) to \(V'\) .
\item
  For every \(a\) belonging to the multiplicative group \(\mathbf{Q}^*\) , let \(\varphi(a)\) be the automorphism of the additive group \(\mathbf{R}\) defined by \(\varphi(a)x = ax\) . Equip \(\mathbf{Q}^*\) with the discrete topology. Let \(\mathbf{G}\) be the topological semi-direct product of \(\mathbf{Q}^*\) and \(\mathbf{R}\) defined by \(\varphi\) (GT, III, §2, No.~10). Show that \(\mathbf{G}\) is locally isomorphic to \(\mathbf{R}\) , but is not unimodular.
\item
  Let \(\mathbf{G}\) be a locally compact group having an open and compact subgroup \(\mathbf{H}\) . Show that for every automorphism \(\varphi\) of \(\mathbf{G}\) , mod \(\varphi \in \mathbf{Q}_+^*\) . (Observe that \(\varphi(\mathbf{H}) \cap \mathbf{H}\) has finite index in \(\mathbf{H}\) and in \(\varphi(\mathbf{H})\) .) Show that the equality \(\Delta_{\mathbf{G}}(s) = 1\) defines an open subgroup of \(\mathbf{G}\) containing \(\mathbf{H}\) .
\item
  Let G be a locally compact group, \(\beta\) a left Haar measure on G, A a subset of G, and B a relatively compact \(\beta\) -integrable subset of G such that \(\beta(\mathrm{B}) > 0\) . Show that if the compact subsets of \(\mathbf{A} \cdot \mathbf{B}\) have bounded measures for \(\beta\) , then \(\mathbf{A}\) is relatively compact. (Imitate the argument of Prop. 2.)
\item
  Let G be a locally compact group, A a dense subset of G, \(\alpha\) a left Haar measure on G, and H an \(\alpha\) -measurable subset of G having the following property: for every \(s \in A\) , \(sH \cap (\mathbb{C}H)\) and \(H \cap (\mathbb{C}sH)\) are locally \(\alpha\) -negligible. Show that either H is locally \(\alpha\) -negligible or its complement is locally \(\alpha\) -negligible. (Show that \(\varphi_{H} \cdot \alpha\) is left invariant.)
\item
  Adopt the notations of the proof of Th. 1 of No.~2. Let \(\beta\) be the unique left Haar measure on G such that \(\beta(f_0) = 1\) . Show that, for every \(f \in \mathcal{K}(\mathbf{G})\) , \(\mathrm{I}_g(f)\) tends to \(\beta(f)\) with respect to \(\mathfrak{B}\) . (Let \(a\) be a cluster point of \(g \mapsto \mathrm{I}_g(f)\) with respect to \(\mathfrak{B}\) . There exists an ultrafilter \(\mathfrak{U}\) finer than \(\mathfrak{B}\) such that \(\mathrm{I}_g(f)\) tends to \(a\) with respect to \(\mathfrak{U}\) . On the other hand, \(\mathrm{I}_g(f)\) tends to \(\beta(f)\) with respect to \(\mathfrak{U}\) .)
\item
  Let G be a locally compact group, \(\mu\) a left Haar measure on G. Show that every \(\mu\) -integrable set A is contained in a countable union of compact sets. (Reduce to the case that A is open, and observe that there exists an open subgroup of G that is a countable union of compact subsets.)
\end{enumerate}

¶ 11) Let G be a nondiscrete locally compact group countable at infinity, β a left Haar measure on G. Show that every compact neighborhood V of e contains a normal subgroup H of G, β-negligible and such that G/H is a Polish locally compact group. (One can proceed as follows: let L be the open subgroup of G generated by V; let \(b_{1}, b_{2}, \ldots\) be representatives of the left cosets with respect to L. Form a decreasing sequence \((\mathrm{V}_{n})\) of symmetric neighborhoods of e such that: 1) \(\mathrm{V}_{n}^{2} \subset \mathrm{V}_{n-1}\) ; 2) \(x\mathrm{V}_{n}x^{-1} \subset \mathrm{V}_{n-1}\) for every \(x \in \mathrm{V}\) ; 3) \(b_{i}\mathrm{V}_{n}b_{i}^{-1} \subset \mathrm{V}_{n-1}\) for \(1 \leqslant i \leqslant n\) ; 4) \(\beta(\mathrm{V}_{n}) \leqslant 1/n\) . Set \(\mathrm{H} = \mathrm{V}_{1} \cap \mathrm{V}_{2} \cap \cdots\) .

Let \((f_{n})\) be a sequence of numerical functions uniformly continuous for the left uniform structure of G. Show that H can be constructed in such a way that, in addition, the \(f_{n}\) are constant on the cosets of H. (In the preceding construction, take \(V_{n}\) to be such that \(|f_{i}(x) - f_{i}(y)| \leqslant 1/n\) for \(x^{-1}y \in V_{n}\) and \(1 \leqslant i \leqslant n\) .)

Let \((g_n)\) be a sequence of \(\beta\) -integrable numerical functions on \(G\) . Show that \(H\) can be constructed in such a way that, in addition, each \(g_n\) is equal almost everywhere to a function that is constant on the cosets of \(H\) .

\begin{enumerate}
\def\labelenumi{\arabic{enumi})}
\setcounter{enumi}{11}
\tightlist
\item
  Let \(\mathbf{K}\) be a nondiscrete locally compact field, \(\mathbf{E}\) a left topological vector space over \(\mathbf{K}\) .
\end{enumerate}

\begin{enumerate}
\def\labelenumi{\alph{enumi})}
\item
  If E is 1-dimensional, then E is isomorphic to \(\mathbf{K}_s\) (imitate the proof of Prop. 2 of TVS, I, §2, No.~2, replacing the absolute value by the function mod \(_K\) ).
\item
  If E is n-dimensional, then E is isomorphic to \(K_{s}^{n}\) (imitate the proof of Th. 2 of TVS, I, §2, No.~3).
\item
  Assume that E is locally compact. Let F be a linear subspace of E of finite dimension n.~For every \(a \in K\) , let \(\text{mod}_{\mathbf{E}}(a)\) be the modulus of \(x \mapsto ax\) in E. Since F is closed in E by b), one can form \(\text{mod}_{\mathbf{E}/\mathbf{F}}(a)\) . Show that \(\text{mod}_{\mathbf{E}}(a) = \text{mod}_{\mathbf{K}}(a)^n \text{mod}_{\mathbf{E}/\mathbf{F}}(a)\) . Show that if \(\text{mod}_{\mathbf{K}}(a) < 1\) and \(F \neq E\) , then \(\text{mod}_{\mathbf{E}/\mathbf{F}}(a) < 1\) . (One has \(a^n \to 0\) as \(n \to +\infty\) ; from this, deduce that \(\text{mod}_{\mathbf{E}/\mathbf{F}}(a^n) \to 0\) by arguing as in Prop. 12 of No.~10.) Deduce from this that \(n \leqslant \log \text{mod}_{\mathbf{E}}(1/a)/\log \text{mod}_{\mathbf{K}}(1/a)\) , hence that E is finite-dimensional.
\end{enumerate}

(We shall see in CA, VI, that the topology of a locally compact field can be defined by an absolute value. The final result of c) will then be a special case of Th. 3 of TVS, I, §2, No.~4.)

¶ 13) Let G be a locally compact group operating continuously on the left in a locally compact Polish space T, β a left Haar measure on G, and ν a positive measure on T quasi-invariant under G.

\begin{enumerate}
\def\labelenumi{\alph{enumi})}
\tightlist
\item
  There exists a function \((s,x)\mapsto \chi (s,x) > 0\) on \(\mathbf{G}\times \mathbf{T}\) , locally \((\beta \otimes \nu)\) -integrable, such that for every \(s\in \mathbf{G}\) , the function \(x\mapsto \chi (s^{-1},x)\) is locally \(\nu\) -integrable and satisfies \(\gamma(s)\nu = \chi(s^{-1},\cdot)\cdot\nu\) . (Show that \((s,x)\mapsto(s,sx)\) transforms \(\beta\otimes\nu\) into an equivalent measure, by making use of criterion 3) of Th. 2 of Ch. V, §5, No.~5; let \(\chi(s^{-1},x)d\beta(s)d\nu(x)\) be this equivalent measure. For \(\varphi\in\mathcal{H}(\mathrm{G})\) and \(\psi\in\mathcal{H}(\mathrm{T})\) ,
\end{enumerate}

\[
\int \varphi (s) d \beta (s) \int \psi (s x) d \nu (x) = \int \varphi (s) d \beta (s) \int \psi (x) \chi (s ^ {- 1}, x) d \nu (x),
\]

whence \(\int \psi(sx) d\nu(x) = \int \psi(x)\chi(s^{-1}, x) d\nu(x)\) except on a locally \(\beta\) -negligible set \(N(\psi)\) . Next make use of Lemma 1 of Ch. VI, §3, No.~1; then modify \(\chi\) on a \((\beta \otimes \nu)\) -negligible set.)

\begin{enumerate}
\def\labelenumi{\alph{enumi})}
\setcounter{enumi}{1}
\tightlist
\item
  Show that for any \(s, t\) in \(\mathbf{G}\) , one has
\end{enumerate}

\[
\chi (s t, x) = \chi (s, t x) \chi (t, x)
\]

except on a \(\nu\) -negligible set of values of \(x\) (make use of the relation \(\pmb{\gamma}(st)\nu = \pmb{\gamma}(s)(\pmb{\gamma}(t)\nu)\) ). c) Show that the function \(\chi\) of \(a\) ) is determined up to a locally \((\beta \otimes \nu)\) -negligible subset of \(\mathbf{G} \times \mathbf{T}\) .

¶ 14) Let T be a locally compact space, U a uniform structure on T for which there exists an entourage \(V_{0}\) such that \(V_{0}(t)\) is compact for all \(t \in T\) (GT, II, §4, Exer. 9). On the other hand, let \(\Gamma\) be a uniformly equicontinuous group of homeomorphisms of T (for the uniform structure U), such that there exists an \(a \in T\) whose orbit under \(\Gamma\) is dense. Show that there exists on T a positive measure \(\neq 0\) invariant under \(\Gamma\) , and that this measure is unique up to a constant factor. One may proceed as follows:

\(1^{\circ}\) As for the existence of the invariant measure, follow the method of No.~2, Th. 1; one proves first that if K is a compact subset of T and U is an open neighborhood of \(a\) , there exists a finite number of elements \(\sigma_{i} \in \Gamma\) such that \(\mathrm{K} \subset \bigcup \sigma_{i}(\mathrm{U})\) .

\(2^{\circ}\) As for uniqueness, observe first that the entourages of \(\mathcal{U}\) that are invariant under every homeomorphism \(\sigma \times \sigma\) of T × T, for \(\sigma \in \Gamma\) , form a fundamental system of entourages \(\mathfrak{S}\) . For every relatively compact set A ⊂ T and every relatively compact set B ⊂ T with nonempty interior, let (A : B) be the smallest number of elements of a covering of A by sets of the form \(\sigma\) B, where \(\sigma \in \Gamma\) ; if C ⊂ T is a third relatively compact set with nonempty interior, then (A : C) ≤ (A : B)(B : C). Let K be a compact subset of T, L ⊃ K a relatively compact open set, V ∈ \(\mathfrak{S}\) a symmetric open entourage contained in V0 such that V(K) ⊂ L, and W ∈ \(\mathfrak{S}\) a symmetric closed entourage contained in V. For every symmetric entourage U ∈ \(\mathfrak{S}\) such that UW ⊂ V and WU ⊂ V, show that, for every positive measure \(\nu\) on T invariant under \(\Gamma\) ,

\begin{enumerate}
\def\labelenumi{(\arabic{enumi})}
\tightlist
\item
\end{enumerate}

\[
\big (\mathrm{W} (a): \mathrm{U} (a) \big) \nu (\mathrm{K}) \leqslant \big (\mathrm{L}: \mathrm{U} (a) \big) \nu \big (\mathrm{V} (a) \big)\tag{2}
\]

\[
\bigl (\mathrm{K}: \mathrm{U} (a) \bigr) \nu \bigl (\mathrm{W} (a) \bigr) \leqslant \bigl (\mathrm{V} (a): \mathrm{U} (a) \bigr) \nu (\mathrm{L}).
\]

For this, one proves that if \((\sigma_i(\mathrm{U}(a)))\) is a covering of L formed by \((\mathrm{L}:\mathrm{U}(a))\) sets, every \(x\in \mathbf{K}\) belongs to at least \((\mathrm{W}(a):\mathrm{U}(a))\) sets of the form \(\sigma_{i}(\mathrm{V}(a))\) , and that every \(y\in \mathrm{L}\) belongs to at most \((\mathrm{V}(a):\mathrm{U}(a))\) sets of the form \(\sigma_{i}(\mathrm{W}(a))\) (note that if \(y\in \sigma_i(\mathrm{W}(a))\) , then \(\sigma_{i}(\mathrm{U}(a))\) is contained in \(\mathrm{V}(x)\) ).

Let \(\mathfrak{F}\) be an ultrafilter finer than the section filter of \(\mathfrak{S}\) ; let \(A_0\) be a nonempty, relatively compact open set in T, and set \(\lambda_U(A) = (A: U(a)) / (A_0: U(a))\) for every symmetric entourage \(U \in \mathfrak{S}\) and every relatively compact subset A of T. Let \(\lambda(A) = \lim_{\mathfrak{F}} \lambda_U(A)\) ; for every compact subset K of T, set \(\lambda'(K) = \inf \lambda(B)\) , where B runs over the set of relatively compact open neighborhoods of K. Deduce from (1) and (2) that

\[
\begin{array}{r} \lambda^ {\prime} \big (\mathrm{W} (a) \big) \nu (\mathrm{K}) \leqslant \lambda (\mathrm{L}) \nu \big (\mathrm{W} (a) \big) \\ \lambda^ {\prime} (\mathrm{K}) \nu \big (\mathrm{W} (a) \big) \leqslant \lambda^ {\prime} \big (\mathrm{W} (a) \big) \nu (\mathrm{L}). \end{array}
\]

Conclude from this that if \(K_{1}, K_{2}\) are two compact subsets of T, and \(L_{1} \supset K_{1}\) , \(L_{2} \supset K_{2}\) are two relatively compact open sets, then

\[
\lambda^ {\prime} (\mathrm{K} _ {1}) \nu (\mathrm{K} _ {2}) \leqslant \lambda (\mathrm{L} _ {2}) \nu (\mathrm{L} _ {1}),
\]

and finally that \(\lambda^{\prime}(\mathrm{K}_{1})\nu (\mathrm{K}_{2}) = \lambda^{\prime}(\mathrm{K}_{2})\nu (\mathrm{K}_{1})\)

¶ 15) Let G be a locally compact group, H a closed subgroup of G, so that G operates continuously on G/H. Assume that H satisfies the following condition: for every neighborhood V of e in G, there exists a neighborhood U of e such that HU ⊂ VH (note that this condition is satisfied for every subgroup H of G if e admits a fundamental system of neighborhoods invariant under every inner automorphism of G). Show that there then exists a nonzero positive measure on G/H invariant under G. (Use the same method as in the proof of Th. 1 of No.~2, on observing, on the one hand, that as U runs over the set of neighborhoods of e in G, the images of the sets HUH under the canonical mapping π : G → G/H form a fundamental system of neighborhoods of π(H); on the other hand, for every neighborhood V of e in G and every compact subset K of G, there exist a neighborhood U of e in G and a compact subset L of G such that every set of the form sHUH that intersects KH is of the form s'HUH with s' ∈ L.)

¶ 16) Let G be a compact abelian group, E a dense subset of G that is stable under the law of composition of G. Assume that, for every bounded numerical function f on E, there is defined a number M(f) having the following property: for any two bounded functions f, g on E, for any elements \(a_{i}, b_{k}\) of E ( \(1 \leqslant i \leqslant m, 1 \leqslant k \leqslant n\) ), and for any real numbers \(\alpha_{i}, \beta_{k}\) ( \(1 \leqslant i \leqslant m, 1 \leqslant k \leqslant n\) ) such that

\[
\sum_ {i} \alpha_ {i} f (x a _ {i}) \leqslant \sum_ {k} \beta_ {k} g (x b _ {k})
\]

for all \(x \in \mathbf{E}\) , one then has the inequality

\[
\left(\sum_ {i} \alpha_ {i}\right) \mathrm{M} (f) \leqslant \left(\sum_ {k} \beta_ {k}\right) \mathrm{M} (g).
\]

Assume in addition that \(\mathbf{M}(1) = 1\) (cf.~TVS, IV, Appendix).

\begin{enumerate}
\def\labelenumi{\alph{enumi})}
\item
  Let \(\mu\) be the normalized Haar measure on G. Show that for every continuous numerical function \(f\) on G, \(\int f d\mu = \mathrm{M}(f|\mathrm{E})\) . (By considering a suitable partition of unity, approximate the constant function equal to \(\int f d\mu\) on E by means of functions of the form \(x \mapsto \sum_{i} \alpha_i f(xa_i)\) , where \(a_i \in \mathrm{E}\) .)
\item
  Deduce from a) that if one sets \(\nu(B) = M(\varphi_B)\) for every subset B of E, then \(\nu(U \cap E) \geqslant \mu(U)\) for every open subset U of G, and \(\nu(F \cap E) \leqslant \mu(F)\) for every closed subset F of G. For every \(\mu\) -quadratic subset P of G (Ch. IV, §5, Exer. 17 d)), one has \(\nu(P \cap E) = \mu(P)\) .
\end{enumerate}

¶ 17) Let T be a locally compact space, S a subset of T equipped with a law of composition \((x,y)\mapsto xy\) that makes it a monoid \(^{1}\) (not necessarily having a priori a neutral element), and which is continuous on S×S when S is equipped with the topology induced by that of T. Assume in addition that every element of S is cancellable \(^{2}\) (A, I, §2, No.~2). Let \(\mu\) be a bounded positive measure on T, concentrated on S (so that S is \(\mu\) -measurable) and of total mass 1. Assume that \(\mu\) is left invariant in the following sense: for every numerical function f defined on S, continuous and bounded (hence \(\mu\) -measurable by Ch. IV, §5, No.~5, Cor. of Prop. 8), one has \(\int f(sx) d\mu(x) = \int f(x)d\mu(x)\) for all \(s \in S\) . It comes to the same to say that \(\mu(sK) = \mu(K)\) for every compact subset K of S.

\begin{enumerate}
\def\labelenumi{\alph{enumi})}
\item
  Consider the product measure \(\mu\otimes\mu\) on T×T; show that there exist compact subsets K of S such that the images of K×K under the two continuous mappings \((x,y)\mapsto(x,xy)\) and \((x,y)\mapsto(xy,x)\) have measure arbitrarily close to 1 (use the Lebesgue--Fubini theorem). Conclude from this that these images have a common point, and deduce therefrom that S has a neutral element (cf.~A, I, §2, Exer. 9).
\item
  Show that S is a compact group. (First prove that for every \(x \in S\) , the inner measure \(\mu_{*}(xS)\) is equal to 1, hence that xS is measurable and that the measure is concentrated on xS; xS is then a monoid to which the result of a) can be applied, which shows that S is a group. Argue as in Prop. 2 to prove that S is compact. Finally, make use of Exer. 21 of GT, III, §4.3)
\end{enumerate}

18) Let X be a locally compact space, G a compact group operating continuously on X, E the orbit of a point of X, and \(\mathcal {F}\) a vector space of continuous numerical functions on X such that, for every function \(f \in \mathcal {F}\) and every \(s \in G\) , one has \(\gamma (s) f \in \mathcal {F}\) ; assume in addition that \(\mathcal {F}\) contains the constant functions on X. Let \(x _ {0} \in \mathrm{X}\) be a point invariant under G and such that \(| f (x _ {0}) | \leqslant \sup _ {y \in E} | f (y) |\) for every function \(f \in \mathcal {F}\) . Show

that one then has \(f(x_0) = \int_{\mathbf{G}} f(s \cdot z) ds\) for every \(z \in \mathbf{E}\) and every \(f \in \mathcal{F}\) . (Show that there exists a positive measure \(\nu\) on \(\mathbf{E}\) , of total mass 1, such that \(f(x_0) = \int_{\mathbf{E}} f(z) d\nu(z)\) for every \(f \in \mathcal{F}\) , and apply the Lebesgue-Fubini theorem.) *The case that \(X = R^n\) , \(G\) is the orthogonal group, and \(x_0 = 0\) ; apply the formula (7) of §2.*

\begin{enumerate}
\def\labelenumi{\arabic{enumi})}
\setcounter{enumi}{18}
\tightlist
\item
  Let X be a compact space, A a normed algebra over R with unity element, and G a compact group; assume that G operates continuously on A and on X and that, for every \(s \in G\) , \(a \mapsto s \cdot a\) is an automorphism of the algebra A such that \(\|s \cdot a\| = \|a\|\) for all \(a \in A\) . A mapping f of X into A is said to be covariant under G if
\end{enumerate}

\[
f (s \cdot x) = s \cdot f (x)
\]

for all \(s \in G\) and \(x \in X\) . Let \(B\) be a subring of \(A^X\) consisting of covariant continuous functions and containing all of the covariant continuous mappings of \(X\) into \(R\) ( \(R\) being identified with a subalgebra of \(A\) ; one observes that for such a mapping \(g\) , \(g(s \cdot x) = g(x)\) for all \(s \in G\) and \(x \in X\) ). Let \(f\) be a covariant continuous mapping of \(X\) into \(A\) , and assume that for every \(y \in X\) , there exists a mapping \(g_y \in B\) such that \(f(y) = g_y(y)\) . Then, for every \(\varepsilon > 0\) , there exists a mapping \(g \in B\) such that \(\|f(x) - g(x)\| \leqslant \varepsilon\) for all \(x \in X\) . (Make use of a suitable continuous partition of unity ( \(\varphi_i\) ) on \(X\) and introduce the functions \(h_i\) such that \(h_i(x) = \int_G \varphi_i(s \cdot x) ds\) .)

¶ 20) Let G be a locally compact group, μ a left Haar measure on G, A and B two subsets of G.

\begin{enumerate}
\def\labelenumi{\alph{enumi})}
\tightlist
\item
  Assume that one of the following two conditions holds:
\end{enumerate}

\(\alpha)\) A is \(\mu\) -integrable;

\(\beta)\ \mu^{*}(A) < +\infty\) and \(B\) is \(\mu\) -measurable.

Show that, in each of these two cases, the function \(f(s) = \mu^{*}(sA \cap B)\) is uniformly continuous on \(G\) for the right uniform structure of \(G\) . (For any two subsets \(M, N\) of \(G\) , one sets

\[
d (\mathrm{M}, \mathrm{N}) = \mu^ {*} \big ((\mathrm{M} \cap \mathbb {C N}) \cup (\mathrm{N} \cap \mathbb {C M}) \big).
\]

First consider the case that A is compact; making use of Ch. IV, §4, No.~6, Th. 4, show that for every \(\varepsilon > 0\) , there exists a neighborhood U of \(e\) in G such that, for every \(s \in G\) and \(t \in U\) ,

\[
d (s \mathrm{A} \cap \mathrm{B}, s t \mathrm{A} \cap \mathrm{B}) \leqslant \varepsilon .
\]

Then apply Exer. 13 of Ch. IV, §5. If B is \(\mu\) -measurable, \(\mu^{*}(\mathbf{A}) < +\infty\) and \((\mathbf{A}_{n})\) is a decreasing sequence of integrable subsets of G containing A such that \(\inf (\mu (\mathbf{A}_n)) = \mu^{*}(\mathbf{A})\) , show that

\[
\mu^ {*} (s \mathrm{A} \cap \mathrm{B}) = \inf \left(\mu^ {*} (s \mathrm{A} _ {n} \cap \mathrm{B})\right)
\]

(Ch. V, 1st edn., §2, No.~2, Lemma 1); \(^{4}\) note on the other hand that \(\mu(A_{n}-A_{n+1})\) tends to 0 with \(1/n\) .)

\begin{enumerate}
\def\labelenumi{\alph{enumi})}
\setcounter{enumi}{1}
\item
  If \(\mathbf{A}^{-1}\) is \(\mu\) -integrable and \(\mu^{*}(\mathbf{B}) < +\infty\) , then the function \(f\) is uniformly continuous for both the right and left uniform structures of \(\mathbf{G}\) ; moreover, one then has \(\int_{\mathbf{G}} f(s) d\mu(s) = \mu(\mathbf{A}^{-1})\mu^{*}(\mathbf{B})\) . (Reduce to the case that \(\mathbf{B}\) is integrable; observe then that \(\mu(s\mathbf{A} \cap \mathbf{B}) = \mu(\mathbf{A} \cap s^{-1}\mathbf{B})\) , that \(\varphi_{s\mathbf{A} \cap \mathbf{B}} = \varphi_{s\mathbf{A}}\varphi_{\mathbf{B}}\) and that \(\varphi_{s\mathbf{A}}(t) = \varphi_{t\mathbf{A}^{-1}}(s)\) .)
\item
  Deduce from \(a\) ) that in the two cases considered, the interiors of AB and BA are not empty if A and B are not negligible. (Cf. Ch. VIII, §4, No.~6, Prop. 17.)
\item
  In the group \(\mathbf{G} = \mathbf{SL}_{2}(\mathbf{R})\) , give an example of a compact set A and a \(\mu\) -measurable set B such that \(f(s) = \mu(sA \cap B)\) is not uniformly continuous for the left uniform structure. (Observe that there exist a sequence \((t_{n})\) of elements of G tending to e and a sequence \((s_{n})\) of elements of G such that the sequence \((s_{n}^{-1}t_{n}s_{n})\) tends to the point at infinity.)
\item
  Give an example of two nonmeasurable sets A, B in a locally compact group G, of finite outer measure and such that the function \(s \mapsto \mu^{*}(sA \cap B)\) is not continuous (cf.~Ch. IV, §4, Exer. 8).
\end{enumerate}

\begin{enumerate}
\def\labelenumi{\arabic{enumi})}
\setcounter{enumi}{20}
\tightlist
\item
  \begin{enumerate}
  \def\labelenumii{\alph{enumii})}
  \tightlist
  \item
    Let G be a locally compact group, \(\mu\) a left Haar measure on G. Show that if S is a stable subset of G such that \(\mu_{*}(S)>0\) , then the interior of S is nonempty (cf.~Exer. 20). In particular, if a subgroup H of G has nonzero inner measure, then H is an open subgroup of G.
  \end{enumerate}
\end{enumerate}

\begin{enumerate}
\def\labelenumi{\alph{enumi})}
\setcounter{enumi}{1}
\item
  If G is compact, then every stable subset S of G such that \(\mu_{*}(S) > 0\) is a compact open subgroup of G (observe that the interior of S is stable, and make use of Exer. 17 b)).
\item
  Assume that G is compact and abelian, written additively, and that there exists in G an element a of infinite order. Show that there exists a stable subset S of G such that \(a \in S\) , \(-a \notin S\) and \(G = S \cup (-S)\) (use Zorn's lemma); deduce from b) that S is not measurable and that \(\mu_{*}(S) = 0\) .
\end{enumerate}

\begin{enumerate}
\def\labelenumi{\arabic{enumi})}
\setcounter{enumi}{21}
\tightlist
\item
  Let G be a locally compact group, \(\mu\) a left Haar measure on G, and A an integrable subset of G such that \(\mu(A) > 0\) . Show that the set H(A) of \(s \in G\) such that \(\mu(A) = \mu(A \cap sA)\) is a compact group. (Observe that H(A) is closed in G, with the help of Exer. 20. To see that H(A) is compact, consider a compact subset B of A such that \(\mu(B) > \mu(A)/2\) and prove that \(H(A) \subset BB^{-1}\) .)
\end{enumerate}

¶ 23) Let G be a locally compact abelian group, written additively, μ a Haar measure on G, and A,B two integrable subsets of G.

\begin{enumerate}
\def\labelenumi{\alph{enumi})}
\tightlist
\item
  For every \(s \in \mathbf{G}\) , let
\end{enumerate}

\[
\mathrm{A} ^ {\prime} = \sigma_ {s} (\mathrm{A}, \mathrm{B}) = \mathrm{A} \cup (\mathrm{B} + s), \quad \mathrm{B} ^ {\prime} = \tau_ {s} (\mathrm{A}, \mathrm{B}) = (\mathrm{A} - s) \cap \mathrm{B}.
\]

Show that \(\mu (\mathrm{A}^{\prime}) + \mu (\mathrm{B}^{\prime}) = \mu (\mathrm{A}) + \mu (\mathrm{B})\) and \(\mathrm{A}' + \mathrm{B}'\subset \mathrm{A} + \mathrm{B}\)

\begin{enumerate}
\def\labelenumi{\alph{enumi})}
\setcounter{enumi}{1}
\tightlist
\item
  Assume that 0 belongs to \(\mathbf{A} \cap \mathbf{B}\) . A pair \((\mathbf{A}', \mathbf{B}')\) of integrable subsets of \(\mathbf{G}\) is said to be derived from (A, B) if there exist a sequence \((s_k)_{1 \leqslant k \leqslant n}\) of elements of \(\mathbf{G}\) and two sequences \((\mathbf{A}_k)_{0 \leqslant k \leqslant n}\) , \((\mathbf{B}_k)_{0 \leqslant k \leqslant n}\) of subsets of \(\mathbf{G}\) such that \(\mathbf{A}_0 = \mathbf{A}\) , \(\mathbf{B}_0 = \mathbf{B}\) , \(\mathbf{A}_k = \sigma_{s_k}(\mathbf{A}_{k-1}, \mathbf{B}_{k-1})\) , \(\mathbf{B}_k = \tau_{s_k}(\mathbf{A}_{k-1}, \mathbf{B}_{k-1})\) for \(1 \leqslant k \leqslant n\) , \(s_k \in \mathbf{A}_{k-1}\) for \(1 \leqslant k \leqslant n\) , and \(\mathbf{A}' = \mathbf{A}_n\) , \(\mathbf{B}' = \mathbf{B}_n\) . Show that there exists a sequence of pairs \((\mathrm{E}_n, \mathrm{F}_n)\) such that \(\mathbf{E}_0 = \mathbf{A}\) , \(\mathbf{F}_0 = \mathbf{B}\) , \((\mathrm{E}_{n+1}, \mathrm{F}_{n+1})\) is derived from \((\mathrm{E}_n, \mathrm{F}_n)\) , and \(\mu((\mathrm{E}_n - s) \cap \mathrm{F}_n) \geqslant \mu(\mathrm{F}_{n+1}) - 2^{-n}\) for every \(n\) and every \(s \in \mathrm{E}_n\) . Set \(\mathbf{E}_{\infty} = \bigcup_{n} \mathbf{E}_n\) , \(\mathbf{F}_{\infty} = \bigcap_{n} \mathbf{F}_n\) . Show that for every \(s \in \mathrm{E}_{\infty}\) ,
\end{enumerate}

\[
\mu \big ((\mathrm{E} _ {\infty} - s) \cap \mathrm{F} _ {\infty} \big) = \mu (\mathrm{F} _ {\infty}).
\]

\begin{enumerate}
\def\labelenumi{\alph{enumi})}
\setcounter{enumi}{2}
\tightlist
\item
  Assume that \(\mu (\mathrm{F}_{\infty}) > 0\) . Show that the function
\end{enumerate}

\[
f (s) = \mu \big ((\mathrm{E} _ {\infty} - s) \cap \mathrm{F} _ {\infty} \big)
\]

can only take on the values 0 and \(\mu(\mathrm{F}_{\infty})\) , and that the set C of \(s \in G\) such that \(f(s) = \mu(\mathrm{F}_{\infty})\) is open and closed, is such that \(\mu(C) = \mu(\mathrm{E}_{\infty})\) , and is the closure of \(\mathrm{E}_{\infty}\) (use Exer. 20 a) and b)). On the other hand let D be the set of \(s \in \mathrm{F}_{\infty}\) such that the intersection of \(\mathrm{F}_{\infty}\) with every neighborhood of \(s\) has measure \(>0\) . Show that \(\mu(D) = \mu(\mathrm{F}_{\infty})\) and that

\[
\mathbf {E} _ {\infty} + \mathbf {D} \subset \mathbf {C};
\]

from this, deduce that D is contained in the subgroup H(C) defined in Exer. 22, and that H(C) is compact and open in G. Finally, show that C + H(C) = C, that \(\mu(C) \geqslant \mu(A) + \mu(B) - \mu(H(C))\) , and that C ⊂ A + B (consider the measure of E∞ ∩ (c - F∞) for every c ∈ C).

\begin{enumerate}
\def\labelenumi{\alph{enumi})}
\setcounter{enumi}{3}
\tightlist
\item
  Deduce from c) that, for two integrable subsets A, B of G: either \(\mu_{*}(A+B)\geqslant\mu(A)+\mu(B)\) ; or else there exists a compact open subgroup H of G such that \(A+B\) contains a coset of H, in which case \(\mu_{*}(A+B)\geqslant\mu(A)+\mu(B)-\mu(H)\) . The case that G is connected.
\end{enumerate}

¶ 24 a) In R, let A (resp. B) be the set of numbers x whose proper dyadic expansion \(x = x_0 + \sum_{i=1}^{\infty} x_i 2^{-i}\) ( \(x_0\) an integer, \(x_i = 0\) or \(x_i = 1\) for \(i \geqslant 1\) ) is such that \(x_{i}=0\) for i even and \textgreater0 (resp. i odd). Show that, for Lebesgue measure, A and B have measure zero but \(A+B=R\) .

\begin{enumerate}
\def\labelenumi{\alph{enumi})}
\setcounter{enumi}{1}
\item
  Deduce from a) that there exists a basis H of R (over Q) contained in A ∪ B, hence of measure zero. The set P₁ of numbers rh, where r ∈ Q and h ∈ H, is also of measure zero.
\item
  Denote by \(P_{n}\) the set of real numbers at most n of whose coordinates relative to the basis H are nonzero. Show that if \(P_{n}\) is negligible and \(P_{n+1}\) is measurable, then \(P_{n+1}\) is negligible. (Let \(h_{0} \in H\) ; show first that the set S of \(x \in P_{n+1}\) whose coordinate relative to \(h_{0}\) is \(\neq 0\) , is negligible. Using Exer. 20, show that if \(P_{n+1}\) were not negligible, there would exist two distinct points \(x', x''\) of \(P_{n+1} \cap CS\) such that \((x' - x'') / h_{0}\) is rational, and deduce from this a contradiction.)
\item
  Deduce from a) and b) that there exist in R two negligible sets C, D such that C + D is not measurable.
\end{enumerate}

¶ 25) a) Let f be a positive numerical function defined on R, integrable (for the Lebesgue measure μ on R), bounded, and with compact support. Let γ = sup f(t).

For every \(w \in \mathbf{R}\) , denote by \(\mathrm{U}_f(w)\) the set of \(t \in \mathbf{R}\) such that \(f(t) \geqslant w\) ; set \(\nu_f(w) = \mu^*(\mathrm{U}_f(w))\) . Show that for every \(\alpha > 1\) ,

\[
\int_ {- \infty} ^ {+ \infty} f ^ {\alpha} (t) d t = \int_ {0} ^ {\gamma} \nu_ {f} (w) \alpha w ^ {\alpha - 1} d w.
\]

\begin{enumerate}
\def\labelenumi{\alph{enumi})}
\setcounter{enumi}{1}
\tightlist
\item
  Let g be a second numerical function satisfying the same conditions as f, and set \(\delta = \sup_{t \in \mathbf{R}} g(t)\) . Let h be the function defined on \(R^{2}\) by \(h(u,v) = f(u) + g(v)\) if \(f(u)g(v) \neq 0\) , and \(h(u,v) = 0\) otherwise; finally, set \(k(t) = \sup_{u+v=t} h(u,v)\) , so that k is positive, integrable, bounded, and has compact support. Show that for every \(\alpha > 1\) ,
\end{enumerate}

\[
\int_ {- \infty} ^ {+ \infty} k ^ {\alpha} (t) d t \geqslant (\gamma + \delta) ^ {\alpha} \left(\frac {1}{\gamma^ {\alpha}} \int_ {- \infty} ^ {+ \infty} f ^ {\alpha} (t) d t + \frac {1}{\delta^ {\alpha}} \int_ {- \infty} ^ {+ \infty} g ^ {\alpha} (t) d t\right).
\]

(Observe that for \(0 < w < 1\) , one has \(\mathrm{U}_k(\gamma w + \delta w) \supset \mathrm{U}_f(\gamma w) + \mathrm{U}_g(\delta w)\) , and use \(a\) and Exer. 23 d).)

\begin{enumerate}
\def\labelenumi{\alph{enumi})}
\setcounter{enumi}{2}
\tightlist
\item
  Let \(\mu_{n}\) be Lebesgue measure on \(\mathbf{R}^{n}\) , A and B two integrable subsets of \(\mathbf{R}^{n}\) . Show that (Brunn-Minkowski inequality)
\end{enumerate}

\[
\left(\left(\mu_ {n}\right) _ {*} (\mathrm{A} + \mathrm{B})\right) ^ {1 / n} \geqslant \left(\mu_ {n} (\mathrm{A})\right) ^ {1 / n} + \left(\mu_ {n} (\mathrm{B})\right) ^ {1 / n}.
\]

(Reduce to the case that A and B are compact. Then argue by induction on n using Exer. 23 d), the Lebesgue--Fubini theorem, the inequality proved in b), as well as Hölder's inequality.)

¶ 26) Let G be a locally compact group, μ a left Haar measure on G, k an integer ≥1, and A an integrable set. Show that for every ε \textgreater{} 0, there exists a neighborhood U of e in G having the following property: for every finite subset S of k elements in U, the set of s ∈ G such that sS ⊂ A has measure ≥(1 - ε)μ(A). (Reduce to the case that A is compact and take U to be such that

\[
\mu (\mathrm{AU}) \leqslant \left(1 - \frac {\varepsilon}{k - 1}\right) \mu (\mathrm{A}).
\]

For every finite subset H of G, set \(p(\mathrm{H}) = \operatorname{Card}(\mathrm{H})\) and \(q(\mathrm{H}) = 1\) if \(\mathrm{H} \neq \emptyset\) , \(q(\mathrm{H}) = 0\) if \(\mathrm{H} = \emptyset\) ; by evaluating the integrals

\[
\int_ {\mathbf {G}} p (\mathrm{A} \cap s \mathrm{S}) d s \quad \text { and } \quad \int_ {\mathbf {G}} q (\mathrm{A} \cap s \mathrm{S}) d s  ,
\]

show that if, for \(h \leqslant k\) , \(\mathbf{M}_h\) denotes the set of \(s \in G\) for which \(\operatorname{Card}(\mathbf{A} \cap s\mathbf{S}) = h\) , then

\[
\sum_ {h = 1} ^ {k} h \mu (\mathrm{M} _ {h}) = k \cdot \mu (\mathrm{A}) \quad \text { and } \quad \sum_ {h = 1} ^ {k} \mu (\mathrm{M} _ {h}) \leqslant \mu (\mathrm{AU}).
\]

From this, conclude that \(\sum_{h=1}^{k-1} h\mu(M_h) \leqslant k\varepsilon \cdot \mu(A)\) .

\begin{enumerate}
\def\labelenumi{\arabic{enumi})}
\setcounter{enumi}{26}
\tightlist
\item
  Let \(\mathbf{G}\) be a group operating on a set \(\mathbf{X}\) . A subset \(\mathbf{P}\) (resp. \(\mathbf{C}\) ) of \(\mathbf{X}\) is said to be a G-filling (resp. G-covering) if, for every \(s \neq e\) in \(\mathbf{G}\) , one has \(s\mathbf{P} \cap \mathbf{P} = \emptyset\) (resp. if \(\mathbf{X} = \bigcup_{s \in \mathbf{G}} s\mathbf{C}\) ). One calls G-paving a subset \(\mathbf{P}\) that is both a G-filling and a G-covering.
\end{enumerate}

\begin{enumerate}
\def\labelenumi{\alph{enumi})}
\item
  Assume that X is locally compact, that G is countable, operates continuously in X, and that there exists a nonzero positive measure \(\mu\) on X invariant under G. Let P and C be a G-filling and a G-covering that are \(\mu\) -integrable. Show that \(\mu(C) \geqslant \mu(P)\) . (Observe that \(\mu(C) \geqslant \sum_{s \in G} \mu(C \cap sP)\) .)
\item
  Assume in addition that there exists on X a uniform structure compatible with the topology of X and admitting a fundamental system S of open entourages invariant under G. On the other hand, one denotes by \(\Delta(G)\) the infimum of the numbers \(\mu(C)\) over all the integrable G-coverings C of X. Let V be an entourage belonging to S, and let \(a \in X\) be such that \(\mu(V(a)) > \Delta(G)\) ; show that there exists an \(s \in G\) such that \(s \neq e\) and \((a, sa) \in \overline{V} \circ V\) .
\item
  Assume that X is a locally compact group, \(\mu\) a left Haar measure, and G a countable subgroup of X, operating by left translation. With the same definition of \(\Delta(G)\) as in \(b\) ), show that if A is an integrable subset of X such that \(\mu(A) > \Delta(G)\) , then there exists \(s \in G \cap AA^{-1}\) such that \(s \neq e\) . Special case that \(X = R^n\) and G is a discrete subgroup of rank \(n\) in \(\mathbf{R}^n: \Delta(G)\) is then equal to the absolute value of the determinant of a basis of G over Z with respect to the canonical basis of \(\mathbf{R}^n\) ; if A is a symmetric closed convex set with nonempty interior \(^5\) in \(\mathbf{R}^n\) such that \(\mu(A) \geqslant 2^n \Delta(G)\) , show that there exists a point of \(A \cap G\) distinct from 0 (Minkowski's theorem).
\item
  With hypotheses as in a), let \(f\) be a function \(\geqslant 0\) and \(\mu\) -integrable on \(X\) . Show that there exist two points \(a, b\) of \(X\) such that \(\mu(C) \sum_{s \in G} f(sa) \geqslant \int_X f(x) d\mu(x)\) and \(\mu (\mathrm{P})\sum_{s\in \mathbb{G}}f(sb)\leqslant \int_{\mathbb{X}}f(x)d\mu (x).\) (Observe that if \(g\) is an integrable function \(\geqslant 0\) , E a
\end{enumerate}

\(\mu\) -integrable set in \(\mathbf{X}\) , then there exists a \(c \in \mathbf{E}\) such that \(\int_{\mathbf{E}} g(x) d\mu(x) \leqslant g(c)\mu(\mathbf{E})\) , and a \(c' \in \mathbf{E}\) such that \(\int_{\mathbf{E}} g(x) d\mu(x) \geqslant g(c')\mu(\mathbf{E})\) .

¶ 28) With hypotheses as in Exer. 27 a), assume in addition that there exists a μ-integrable G-paving F.

\begin{enumerate}
\def\labelenumi{\alph{enumi})}
\tightlist
\item
  For every \(\mu\) -integrable numerical function \(f \geqslant 0\) defined on X, set \(\widetilde{f}(x) = \sum_{s \in G} f(sx)\) , so that \(\int_{F} \widetilde{f}(x) d\mu(x) = \int_{X} f(x) d\mu(x)\) (§2, No.~10, Prop. 15). Let \((f_i)_{1 \leqslant i \leqslant n}\) be a family of numerical functions \(\geqslant 0\) , \(\mu\) -integrable on X: for \(i \neq j\) set
\end{enumerate}

\[
m _ {i j} = \int_ {\mathbf {F}} \widetilde {f} _ {i} (x) \widetilde {f} _ {j} (x) d \mu (x),
\]

and let \(c_{i} = \sup_{x\in X}\widetilde{f}_{i}(x)\) . Show that there exists at least one pair of distinct indices \((i,j)\) such that

\[
m _ {i j} \geqslant \frac {1}{n (n - 1) \mu (\mathrm{F})} \left(\left(\sum_ {i = 1} ^ {n} \int_ {\mathrm{X}} f _ {i} (x) d \mu (x)\right) ^ {2} - \mu (\mathrm{F}) \sum_ {i = 1} ^ {n} c _ {i} \int_ {\mathrm{X}} f _ {i} (x) d x\right).
\]

(Bound the sum \(\sum_{i\neq j}m_{ij}\) from below, using the Cauchy-Schwarz inequality.) From this, deduce that if \((\mathrm{A}_i)_{1\leqslant i\leqslant n}\) is a finite family of \(\mu\) -integrable G-fillings, and if

\[
\sum_ {i = 1} ^ {n} \mu (\mathrm{A} _ {i}) > \mu (\mathrm{F}),
\]

then there exist a pair \((i,j)\) of distinct indices and an \(s\in G\) such that \(\mu (\mathrm{A}_is\cap \mathrm{A}_j) > 0\) b) If \(\mathbf{G}_0\) is a subgroup of \(G\) of finite index \((G:\mathbf{G}_0) = h\) , and if \((s_1,\ldots ,s_h)\) is a system of representatives of the right cosets of \(\mathbf{G}_0\) in \(G\) , then \(\mathrm{F_0} = \bigcup_{1\leqslant i\leqslant h}s_i\mathrm{F}\) is a \(G_{0}\)-paving.

\begin{enumerate}
\def\labelenumi{\alph{enumi})}
\setcounter{enumi}{2}
\tightlist
\item
  In \(\mathbf{X} = \mathbf{R}^n\) , let \(\mathbf{A}\) be a symmetric closed convex set with nonempty interior; let \(u_i : (x_j) \mapsto \sum_{j=1}^{n} c_{ij} x_j\) be \(m\) linear forms on \(\mathbf{X}\) , with integer coefficients \(c_{ij} (m < n)\) .
\end{enumerate}

Show that for every integer \(p \geqslant 1\) and every number \(r \geqslant 0\) such that \(\mu(A)r^n \geqslant 2^n p^m\) , there exists a point \(x \in rA \cap Z^n\) distinct from 0 and such that \(u_i(x) \equiv 0 \pmod{p}\) for \(1 \leqslant i \leqslant m\) (apply Minkowski's theorem (Exer. 27 c)) to the subgroup \(G_0\) of \(Z^n\) formed by the \(z \in Z^n\) such that \(u_i(z) \equiv 0 \pmod{p}\) for \(1 \leqslant i \leqslant m\) , and make use of \(b\) ). Special case that \(n = 2\) , \(m = 1\) and A is defined by \(|x_1| \leqslant 1\) , \(|x_2| \leqslant 1\) (Thue's theorem).

¶ 29) a) Let p be a prime number; there exist two integers a, b such that \(a^{2}+b^{2}+1 \equiv 0 \pmod{p}\) (Alg., Ch. V, 1st edn., §11, No.~5, Cor. of Th. 3). Show that there exist integers \(x_{1}, x_{2}, x_{3}, x_{4}\) not all zero such that \(ax_{1} + bx_{2} \equiv x_{3} \pmod{p}\) , \(bx_{1} - ax_{2} \equiv x_{4} \pmod{p}\) and

\[
y = x _ {1} ^ {2} + x _ {2} ^ {2} + x _ {3} ^ {2} + x _ {4} ^ {2} \leqslant \sqrt {2} p
\]

(same method as in Exer. 28 c)). Show that \(y\) is divisible by \(p\) , and deduce therefrom that \(y = p\) .

\begin{enumerate}
\def\labelenumi{\alph{enumi})}
\setcounter{enumi}{1}
\tightlist
\item
  Deduce from a) that every integer \(n \geqslant 0\) is the sum of four squares (Lagrange's theorem; make use of Alg., Ch. IV, 1st edn., §2, Exer. 11).
\end{enumerate}

\begin{enumerate}
\def\labelenumi{\arabic{enumi})}
\setcounter{enumi}{29}
\tightlist
\item
  Let G be a locally compact group, \(\mu\) a left Haar measure on G, and \(\nu\) a bounded measure on G. Assume that the mapping \(s \mapsto \gamma(s)\nu\) of G into the Banach space \(\mathcal{M}^{1}(G)\) is continuous. Show that the measure \(\nu\) has base \(\mu\) . (Let A be a \(\mu\) -negligible compact subset of G. Arguing as for Prop. 11, show that \(\nu(xA)=0\) for x running over a dense subset of G. From this, deduce that \(\nu(A)=0\) .)
\end{enumerate}

Conversely, if \(\nu\) has base \(\mu\) , then the mapping \(s \mapsto \pmb{\gamma}(s)\nu\) of G into \(\mathcal{M}^1(\mathrm{G})\) is continuous: cf.~Ch. VIII, §2, No.~5, Prop. 8.
% semantic-fragments: legacy-input-seam
\relax{}
